\section{Introduction}

The deployment of autonomous multi-agent systems in consequential domains---finance, infrastructure, healthcare, legal reasoning---has exposed a structural gap in the architecture of modern AI systems. These systems can plan, delegate, and execute actions across distributed reasoning trees, yet they consistently lack a mandatory mechanism that compels unresolved exceptions to reach a human principal. This gap is not an oversight; it is a consequence of treating human oversight as a policy choice rather than an architectural requirement. We argue that this distinction determines whether a system is accountable in principle or accountable in fact.

\subsection{The Accountability Gap in Autonomous Multi-Agent Systems}

Modern agentic systems derive their power from compositional delegation: a root agent spawns child agents, which spawn further agents, forming a hierarchical reasoning and execution tree. Each node in this tree is authorized to take consequential actions within its provisioned scope. When a node encounters a condition outside that scope---an unresolved exception---the default behavior in most deployed systems is machine-level suppression: the exception is logged, logged with a recommendation, or suppressed entirely, with the agent proceeding under implicit authorization it was never explicitly granted. The consequence is a class of failures that are architecturally invisible. They are not invisible because the system lacks awareness of them; they are invisible because the system's design grants the machine actor discretion over whether and when a human is informed.

This failure mode is not hypothetical. Systems without mandatory bail-out mechanisms exhibit a consistent pattern: exceptions are silently contained at the lowest machine level capable of containing them, regardless of severity or downstream consequence. The human principal---the only actor with genuine authority to authorize exception recovery---is never notified because the machine actor has determined, on its own terms, that notification was unnecessary. The human principal's authority over exception recovery is therefore nominal, not structural. The machine actor retains real control.

\subsection{Failure Taxonomy Without Mandatory Escalation}

In the absence of an architecturally enforced bail-out mechanism, exception handling in autonomous systems degrades along three distinct failure modes.

The first is \textbf{suppression drift}. An agent encounters an unresolved exception, determines it can proceed without escalating, and does so on the basis of its own risk assessment rather than any external authorization. The exception is not resolved; it is deferred. The system state diverges from the authorized state without human awareness. This drift is invisible precisely because the agent retains discretion over whether to surface the exception.

The second is \textbf{escalation deadlock}. An agent attempts to escalate but finds no defined escalation path, no reachable human principal, or an ambiguous designation of authority. The agent is structurally unable to fulfill an accountability obligation, and the exception state persists indefinitely or until the agent proceeds anyway. The human principal's authority is real but unreachable.

The third is \textbf{log-only compliance}. A system is nominally configured to notify human principals of exceptions, but the notification mechanism is advisory rather than compulsory. The machine actor decides when to surface the flag, which means it can defer or suppress notification precisely when the exception is most consequential. The human is notified when the machine decides the human should be notified.

In each case, the human principal's authority is not abolished; it is structurally subordinated to machine-level discretion. The agent's capacity to contain exceptions exceeds the system's capacity to compel escalation. This asymmetry is not a design flaw; it is the default condition of any system that implements human oversight as a policy rather than an architectural constraint.

\subsection{The Core Insight: Accountability Requires Architectural Fallback}

The central claim of this paper is that accountability in autonomous systems requires an architectural fallback, not merely a policy commitment. A policy-based approach to human oversight states that agents \emph{should} escalate unresolved exceptions. An architectural approach states that agents \emph{cannot proceed} past a certain point without human authorization, and that this constraint is enforced structurally---below the agent's execution level---rather than by the agent itself.

The difference is not one of degree; it is one of kind. Policy-based oversight is a convention that depends on the correct behavior of the very actors it is designed to govern. Architectural enforcement is a compulsion that operates independently of agent volition. A system that relies on policy alone is structurally indistinguishable, at the machine level, from a system with no oversight at all: the agent decides, and the agent's decision is final.

The architectural fallback we propose is the \textbf{Bailout Protocol}, a deterministic state machine embedded in the \fact{} layer of the \bxthree{} Framework. The \fact{} layer is the structural substrate with write access to the immutable forensic ledger and the pre-commit enforcement hooks that govern every consequential action at a node. It is the only layer with authority to reject a directive on safety grounds, and it operates below the agent's execution level. The Bailout Protocol is not invoked by agent code; it is a structural extension of the enforcement substrate that activates when the \bounds{} layer's exception resolution timeout expires without resolution, or when the trigger condition TC is satisfied.

When the Bailout Protocol activates, the system transitions to an \texttt{ESCALATING} state and deterministically propagates the exception along the immutable parent-pointer chain $\alpha(\cdot)$ to the node's designated human anchor $h(n)$. The propagation path bypasses all machine actors in the intervening hierarchy: no agent, parent node, or middleware may interpose, redirect, suppress, or substitute the human anchor in the escalation path. The human principal is the structurally unavoidable terminus of every unresolved exception. This is not a design preference; it is a structural invariant enforced by the \fact{} layer's pre-commit logic.

The three-layer model of the \bxthree{} Framework provides the substrate on which this compulsion is architecturally realized. The \purpose{} layer defines the human anchor and the node's authorized operational scope. The \bounds{} layer defines the operating envelope and houses the Bounds Engine that detects and classifies constraint violations. The \fact{} layer enforces the state machine transitions, records every escalation event immutably, and prevents any machine actor from modifying the escalation path. Without this three-layer separation, the Bailout Protocol would reduce to a discretionary interface whose enforcement depended on the correct behavior of the very actors it is designed to govern. The \bxthree{} Framework is not a container for the Bailout Protocol; it is the mechanism by which the protocol becomes compulsory.

\subsection{Paper Roadmap}

The remainder of this paper is organized as follows. Section~\ref{sec:background} provides the formal problem definition, including the spawn tree model, exception state taxonomy, and the upstream accountability guarantee. Section~\ref{sec:bailout} provides the complete specification of the Bailout Protocol, including the deterministic state machine definition, the trigger condition TC, the escalation algorithm, the bypass property, and timeout handling. Section~\ref{sec:guarantees} establishes the protocol's correctness properties: SAFETY (no consequential action proceeds during unresolved exception), LIVENESS (the protocol reaches the human anchor in finite time), and ACCOUNTABILITY (the escalation path is structurally compulsion rather than voluntary delegation). Section~\ref{sec:limitations} classifies the environmental assumptions that bound the protocol's operational scope, identifies addressable extensions, and charts the formal verification agenda. Section~\ref{sec:related} situates the Bailout Protocol within the landscape of human-in-the-loop AI, HITL control theory, and accountable autonomous systems. Section~\ref{sec:conclusion} concludes with the broader significance of architectural compulsion for the design of accountable AI systems.